% ==========================================
% BAB V IMPLEMENTASI
% ==========================================
\cleardoublepage
\chapter{IMPLEMENTASI \textit{FRONT-END} APLIKASI ITB ULTRA-MARATHON}
\label{chap:implementasi}

Bab ini membahas proses implementasi antarmuka \textit{front-end} aplikasi UltraTrack, sistem \textit{tracker} ITB Ultra-Marathon, berdasarkan rancangan yang telah disusun pada Bab IV. Implementasi mencakup dua platform, yaitu aplikasi \textit{mobile} untuk kebutuhan pelacakan peserta di lapangan dan aplikasi web untuk kebutuhan pemantauan, administrasi, serta akses publik. Seluruh kode dibangun di atas ekosistem JavaScript menggunakan TypeScript sebagai bahasa pemrograman utama.

% ==========================================
\section{Lingkungan Implementasi}
% ==========================================

Lingkungan implementasi yang digunakan dalam pengembangan sistem terbagi menjadi dua, yaitu perangkat keras dan perangkat lunak. Perangkat keras yang digunakan mencakup mesin pengembangan serta perangkat pengujian aplikasi \textit{mobile}. Rincian perangkat keras dapat dilihat pada Tabel \ref{tab:hw}.

\input{tables/hw.tex}

Untuk mendukung pengembangan sistem, perangkat lunak yang digunakan terbagi menjadi dua bagian, yaitu untuk pengembangan aplikasi \textit{mobile} dan pengembangan antarmuka \textit{website}. Rincian perangkat lunak aplikasi \textit{website} dapat dilihat pada Tabel \ref{tab:sw-web}, sedangkan rincian untuk \textit{mobile} dapat dilihat pada Tabel \ref{tab:sw}.

\begin{table}[H]
\centering
\caption{Perangkat Lunak Pengembangan Website}
\label{tab:sw-web}
\begin{tabular}{|p{5cm}|p{8cm}|}
\hline
\textbf{Spesifikasi} & \textbf{Detail} \\ \hline
Bahasa Pemrograman        & TypeScript 6.0 \\ \hline
\textit{Framework} Utama  & React 19 \\ \hline
\textit{Build Tool}       & Vite 8 \\ \hline
Sistem Navigasi           & TanStack Router 1.x \\ \hline
\textit{State Management} & TanStack Query 5.x \\ \hline
\textit{Framework} Styling & Tailwind CSS 4.x \\ \hline
\textit{Library} Peta     & React Leaflet 5.x \\ \hline
Visualisasi Data          & Recharts 2.x \\ \hline
Manajemen Formulir        & React Hook Form 7.x \\ \hline
\textit{Code Editor}      & Visual Studio Code \\ \hline
\textit{Version Control}  & Git \\ \hline
\textit{Package Manager}  & npm \\ \hline
\end{tabular}
\end{table}

\begin{table}[H]
\centering
\caption{Lingkungan Implementasi Perangkat Lunak}
\label{tab:sw}
\begin{tabular}{|p{5cm}|p{8cm}|}
\hline
\textbf{Spesifikasi} & \textbf{Detail} \\ \hline
Bahasa Pemrograman       & TypeScript 5.9 \\ \hline
\textit{Framework} Utama & React Native 0.81.5 \\ \hline
\textit{Platform} Pengembangan & Expo SDK 54 \\ \hline
Sistem Navigasi          & Expo Router 6.0 \\ \hline
\textit{Library} Peta    & React Native Maps 1.20.1 \\ \hline
\textit{Code Editor}     & Visual Studio Code \\ \hline
\textit{Version Control} & Git \\ \hline
\textit{Package Manager} & npm \\ \hline
\textit{Target Platform} & Android (format APK) \\ \hline
\end{tabular}
\end{table}

% ==========================================
\section{Implementasi \textit{Front-End}}
\label{sec:impl-frontend}
% ==========================================

Seluruh \textit{source code} hasil implementasi pada penelitian ini dipublikasikan melalui repositori GitHub dan dapat diakses secara terbuka pada tautan \url{https://github.com/thaliafahira/frontend-itbultramarathon}.

Berikut ini merupakan penjelasan mengenai implementasi \textit{front-end} aplikasi web dan mobile ITB Ultra-Marathon.

% ==========================================
\subsection{Implementasi \textit{Front-End} Aplikasi Web}
\label{sec:impl-web}
% ==========================================

Aplikasi web dibangun menggunakan \textit{framework} React dengan \textit{build tool} Vite. Sistem navigasi menggunakan TanStack Router yang memungkinkan implementasi autentikasi \textit{role-based routing} langsung pada level \textit{router loader}. Antarmuka dibangun dengan menggunakan Tailwind CSS, sedangkan komponen interaktif menggunakan \texttt{radix-ui} untuk memastikan standar aksesibilitas. Untuk peta, digunakan \texttt{react-leaflet} dan \texttt{leaflet}.

Struktur direktori proyek aplikasi web diorganisasikan sebagaimana ditunjukkan pada Listing \ref{lst:struktur-web}.

\begin{lstlisting}[language={},caption={Struktur Direktori Aplikasi Web ITB Ultra-Marathon},label={lst:struktur-web}]
web/
├── index.html                  # Titik masuk utama (Entry point) HTML
├── package.json                # Daftar dependensi library web
├── vite.config.ts              # Konfigurasi Vite bundler
└── src/
    ├── main.tsx                # Entry point React
    ├── router.tsx              # Konfigurasi TanStack Router
    ├── styles.css              # File CSS global (Tailwind)
    ├── components/             # Komponen UI (Reusable)
    │   ├── layout/             # Komponen kerangka halaman (Sidebar, Header)
    │   ├── race/               # Komponen spesifik lomba (AlertCard, StatusBadge)
    │   ├── ui/                 # Komponen dasar (Button, Dialog, Input)
    │   └── RaceMap.tsx         # Komponen utama pe-render peta Leaflet
    ├── data/                   # Data statis (Dummy Data)
    │   ├── captains.ts
    │   ├── race.ts
    │   └── teams.ts
    ├── features/               # Pengelompokan fitur berdasarkan peran
    │   ├── captain/            # Tampilan Ketua Tim
    │   ├── committee/          # Tampilan Panitia
    │   ├── public/             # Tampilan Publik
    │   ├── spectator/          # Tampilan Penonton Individu
    │   └── supporter/          # Tampilan Supporter
    ├── hooks/                  # Custom React Hooks
    │   ├── use-runner-snapshot.ts
    │   └── use-team-live.ts
    ├── lib/                    # Fungsi utilitas
    │   ├── api.ts
    │   ├── auth/
    │   ├── format/
    │   ├── geo/
    │   └── storage/
    ├── routes/                 # Routing halaman
    │   ├── __root.tsx
    │   ├── index.tsx
    │   ├── admin.*.tsx
    │   ├── captain.*.tsx
    │   ├── team.*.tsx
    │   └── view.*.tsx
    └── store/                  # Manajemen state global
        └── race-store.ts
\end{lstlisting}

% ==========================================
\subsection{Implementasi \textit{Front-End} Aplikasi \textit{Mobile}}
\label{sec:impl-mobile}
% ==========================================

Aplikasi \textit{mobile} dibangun menggunakan \textit{framework} React Native dengan platform Expo SDK. Arsitektur navigasi menggunakan Expo Router yang berbasis \textit{file-based routing}. Untuk kebutuhan peta, implementasi menggunakan \texttt{react-native-webview}. \textit{State} aplikasi dikelola oleh \texttt{zustand} dan data sesi disimpan secara persisten menggunakan \texttt{@react-native-async-storage/async-storage}.

Struktur direktori proyek aplikasi \textit{mobile} dapat dilihat pada Listing \ref{lst:struktur-mobile}.

\begin{lstlisting}[language={},caption={Struktur Direktori Aplikasi Mobile ITB Ultra-Marathon},label={lst:struktur-mobile}]
mobile/
├── app.json                    # Konfigurasi aplikasi Expo
├── package.json                # Daftar dependensi library mobile
├── babel.config.js             # Konfigurasi Babel
├── src/
│   ├── components/             # Komponen visual
│   │   ├── atoms/
│   │   ├── molecules/
│   │   └── organisms/
│   ├── constants/              # Variabel konstan
│   │   └── colors.ts
│   ├── data/                   # Data statis
│   │   ├── race.ts
│   │   └── teams.ts
│   ├── hooks/
│   │   └── use-geo-status.ts
│   ├── lib/
│   │   └── api.ts
│   └── store/
│       └── session.ts
└── app/                        # Sistem routing Expo Router
    ├── _layout.tsx
    ├── index.tsx
    ├── invite/
    │   └── index.tsx
    ├── (auth)/                 #login untuk proses development
    │   └── login.tsx
    └── (pelari)/
        ├── _layout.tsx
        ├── home.tsx
        ├── emergency.tsx
        ├── manage.tsx
        └── menu.tsx
\end{lstlisting}

Aplikasi \textit{mobile} yang dikembangkan pada penelitian ini didistribusikan dalam format APK (\textit{Android Package}) untuk platform Android. Berkas APK hasil implementasi dapat diunduh melalui tautan \url{https://bit.ly/3T0IXSE}.

% ==========================================
\subsection{\textit{Deployment} dan Tautan Akses}
\label{sec:impl-deployment}
% ==========================================

Aplikasi web telah di-\textit{deploy} menggunakan layanan Netlify sehingga dapat diakses secara daring melalui tautan \url{https://itbultramarathon.netlify.app/}. Deployment ini memungkinkan aplikasi digunakan dan diuji secara langsung melalui peramban web tanpa memerlukan proses instalasi tambahan.
Setiap peran pengguna mengakses antarmuka melalui tautan yang berbeda sesuai dengan hak aksesnya masing-masing, sebagaimana dirinci pada Tabel \ref{tab:deployment-links}.

\begin{table}[H]
    \centering
    \caption{Tautan Akses Aplikasi Web Berdasarkan Peran}
    \label{tab:deployment-links}
    \begin{tabular}{|l|l|p{6cm}|}
        \hline
        \textbf{No} & \textbf{Peran} & \textbf{Tautan Akses} \\ \hline
        1 & Panitia & \url{https://[domain]/admin} \\ \hline
        2 & Ketua Tim & \url{https://[domain]/captain} \\ \hline
        3 & Supporter & \url{https://[domain]/team} \\ \hline
        4 & Penonton private & \url{https://[domain]/view} \\ \hline
        5 & Penonton Publik & \url{https://[domain]/} \\ \hline
    \end{tabular}
\end{table}
